\documentclass{article}
\usepackage[utf8]{inputenc}
\usepackage[T1]{fontenc}
\usepackage{geometry}
\usepackage{graphicx}
\usepackage{float}
\usepackage{hyperref}
\usepackage[english,indonesian]{babel}

\title{Optimasi Pemrosesan Template LaTeX Menggunakan Pendekatan Berbasis AI}
\author{Academic Support Agent}
\date{\today}

\begin{document}
\maketitle

\begin{abstract}

\end{abstract}

\chapter{PENDAHULUAN}


\chapter{TINJAUAN PUSTAKA}


% Original content (commented out):
% 
%% Preâmbulo (configurações, pacotes e tudo mais)
\input{estrutura/preambulo}

%% Início do Documento
\begin{document}

  %% Documento será escrito em Português do Brasil
  \selectlanguage{brazil}

  %% Elementos pré-textuais: capa, folha de rosto, dedicatória, listas, sumário, etc.
  \input{estrutura/pre-textual.tex}

  %% Indica que a partir daqui ficarão os elementos textuais (TCC em si)
  \textual

  %% Inclui os capítulos do TCC (parte textual)
  \input{estrutura/textual}

  %% Finalizações para o PDF
  \bookmarksetup{startatroot}

  %% Elementos pós-textuais: Referências, Anexos, etc.
  \include{estrutura/pos-textual}

\end{document}

\end{document}